% HD-N16. GPT-6 Astra Pro, /root/imprimitive; inherited effort.
% 21 November 2026. Exact cone quotient and universal Ext surplus.
\subsection{HD-N16: a trace kernel for every six-package graph cone}
\label{sec:HD-N16}

\paragraph{Statement and precise transport scope.}
For every integer $d\geq1$, the six-package objects $E_\delta$ defined
below are actual perfect complexes on $E_i^8$, have no negative
self-Ext, and have whole character
\[
             \operatorname{ch}(E_\delta)=e^\theta+96w_d,
                 \qquad 0\ne w_d\in W_{\mathbf Q(\sqrt{-d})}.
\]
For every actual extension class $\delta$, with arbitrary entries
between all six packages and the ample line summand, their degree-two
categorical trace has a kernel of dimension at least
$70\,732\,408$. In particular none is semiregular. Their original
point stabilizer is trivial. The nonzero trace kernel persists under
actual Fourier--Mukai equivalences and under finite \emph{\'etale}
isogeny pullback equipped with the canonical inherited deck action,
including restriction of that action to subgroups and its conjugation
by those equivalences. This assertion concerns these specified
actions; the calculation does not assert an invariant-dimension bound
for a different coherent linearization newly chosen after isogeny.

\paragraph{1. The actual six packages.}
Use the CM eightfold, integral graph classes and field action of
HD-G21. Explicitly, $A=X\times X$, $X=E_i^4$, $M=\mathcal O_{E_i}(0)$,
and
\[
 L=\boxtimes_{j=1}^4 L_j,\quad
 L_j=p_{1j}^*M^d\otimes p_{2j}^*M,\quad
 \theta_j=c_1(L_j)=d x_jy_j+u_jv_j,\quad \theta=\sum_j\theta_j.
\]
The actual endomorphism $J_d(z,z')=(z',-dz)$ satisfies
$J_d^2=-d$, has holomorphic multiplicities $(4,4)$, and is
skew-adjoint for this polarization. The explicit hyperbolic
isometry of HD-S21 identifies its rational Hermitian space with
$U^4$. No new field action is chosen in this construction.

For $\epsilon\in\{-1,0,1\}^4$, let $B_\epsilon$ be the graph of
$\operatorname{diag}(\epsilon)$ and put
\[
 f_j(z)=(u_j-zx_j)(v_j-zy_j),\quad
 F(s)=\prod_j f_j(s),\quad s^2=-d,\quad
 w_d=\frac{F(s)-F(-s)}{2s}.
\]
The graph's actual kernel map gives
$[B_\epsilon]=\prod_j f_j(\epsilon_j)$; its trivial tangent and
normal bundles give
$\operatorname{ch}(\mathcal O_{B_\epsilon})=[B_\epsilon]$ by
Grothendieck--Riemann--Roch. As checked in HD-G21, $F(s)$ and $F(-s)$
are the two field determinant volumes, and the displayed odd part
is nonzero and integral. The one-pair identity
\[
 f_j(s)=(d+1)f_j(0)+\frac{s-d}{2}f_j(1)
                              +\frac{-s-d}{2}f_j(-1)
\]
gives the exact integral identity
\begin{equation}
 16w_d=\sum_\epsilon c_\epsilon[B_\epsilon],\qquad
 c_\epsilon=2^{4-k}(d+1)^{4-k}S_\epsilon,
 \quad R_\epsilon+S_\epsilon s
        =\prod_{\epsilon_j\ne0}(-d+\epsilon_js),       \tag{N16.1}
\end{equation}
where $k$ counts the nonzero coordinates. In particular
$c_{-\epsilon}=-c_\epsilon$ and
$c_{e_j}=8(d+1)^3>0$.

For each two-element subset $J\subset\{1,2,3,4\}$, define the
actual line bundle and sheaves
\[
 M_J=\bigotimes_{j\notin J}L_j,\quad
 F_{J,\epsilon}=M_J\otimes\mathcal O_{B_\epsilon},\quad
 P_0=\bigoplus_{J,\ c_\epsilon>0}F_{J,\epsilon}^{\,c_\epsilon},
 \qquad
 N=\bigoplus_{J,\ c_\epsilon<0}F_{J,\epsilon}^{\,-c_\epsilon}.
\]
Zero coefficients are omitted. Put $P=L\oplus P_0$ and, for every
$\delta\in\operatorname{Ext}^2(P,N)$, form
\begin{equation}
 N[1]\xrightarrow{i}E_\delta\xrightarrow{q}P
                  \xrightarrow{\delta}N[2],
 \qquad E_\delta=\operatorname{Cone}(\delta)[-1].       \tag{N16.2}
\end{equation}
Multiplication in the four-dimensional exterior algebra of pair $j$
gives $\theta_j f_j(s)=(d+s^2)x_jy_ju_jv_j=0$ and the same
identity for $-s$. Thus $\theta_jw_d=0$ and
$e^{c_1(M_J)}w_d=w_d$. Adding the six copies of (N16.1) proves
the asserted whole character of (N16.2), including all degrees.

Positive and negative slopes are disjoint sets. A map between
line bundles supported on distinct graphs would have image on their
proper intersection, whereas the target is torsion-free on its
irreducible graph. It is therefore zero. Moreover
\begin{equation}
 L^{-1}M_J|_{B_\epsilon}
       =\boxtimes_{j\in J}M^{-(d+\epsilon_j^2)}
          \boxtimes_{j\notin J}\mathcal O_{E_i} .       \tag{N16.3}
\end{equation}
There are exactly two negative factors, so its cohomology occurs
only in degrees $2,3,4$. This proves
$\operatorname{Hom}(P,N)=0$. Also
$\operatorname{Hom}(N,P)=0$: its graph-to-graph terms vanish as
above, and a map from a torsion sheaf to $L$ is zero. All sheaves
are perfect on the smooth variety. The two Hom exact sequences of
(N16.2), exactly as in HD-G21, identify
\[
 \operatorname{Ext}^a(E_\delta,E_\delta)
       \simeq\operatorname{Ext}^{a+1}(P,N)\qquad(a<0).
\]
For $a=-1$ this is the zero Hom space just proved; for smaller $a$
it is negative Ext between sheaves. Thus these are actual gluable
perfect complexes. No restriction on the entries of $\delta$ has
been used.

\paragraph{2. The exact upper-triangular trace-kernel quotient.}
Define the actual linear map
\begin{equation}
 j_\delta:\operatorname{Ext}^3(P,N)\longrightarrow
                  \operatorname{Ext}^2(E_\delta,E_\delta),
 \qquad u\longmapsto i[2]\,u\,q.                       \tag{N16.4}
\end{equation}
Its kernel is exactly
\begin{equation}
 \ker j_\delta=
 \{a[2]\delta+\delta[1]b:
       a\in\operatorname{Ext}^1(N,N),\
       b\in\operatorname{Ext}^1(P,P)\}.                \tag{N16.5}
\end{equation}
Here every product is literal composition with its displayed shift.
To prove this, suppose $i[2]uq=0$. Applying
$\operatorname{Hom}(E_\delta,-)$ to the shifted triangle gives
$uq=\delta[1]a_0$ for a map $a_0:E_\delta\to P[1]$.
Applying $\operatorname{Hom}(-,P[1])$ to the original triangle
identifies this map with $bq$ for
$b\in\operatorname{Ext}^1(P,P)$: the two adjacent groups are
$\operatorname{Ext}^{-1}(N,P)$ and $\operatorname{Hom}(N,P)$,
both zero. Hence $(u-\delta[1]b)q=0$. Applying
$\operatorname{Hom}(-,N[3])$ now gives
$u-\delta[1]b=a[2]\delta$ for
$a\in\operatorname{Ext}^1(N,N)$. Conversely both kinds of terms
map to zero because consecutive maps of (N16.2) compose to zero.
This proves (N16.5), including both inclusions.

Every value of (N16.4) has zero categorical degree-two trace.
Indeed cyclicity of the trace, with $q$ of degree zero, gives
\begin{equation}
 \sigma^2_{E_\delta}(i[2]uq)
       =\sigma^2_P(q[2]i[2]u)=0.                     \tag{N16.6}
\end{equation}
This is the same categorical trace used in Perry's definition.
Under the abelian HKR identification it is the full Atiyah trace
$\sum_p\operatorname{tr}(\xi\operatorname{At}^p/p!)$.
Formula (N16.6) also follows separately in each HKR component by
naturality of the Atiyah class and ordinary graded cyclic trace.
The published comparison is Buchweitz--Flenner,
\emph{The global decomposition theorem for Hochschild (co-)homology
of singular spaces via the Atiyah--Chern character},
Adv. Math. \textbf{217} (2008), 243--281,
Theorem6.4.2, with the trace construction in Proposition6.2.1 and
Corollary6.3.2. These are the identical trace maps already used
and checked in HD-N02 and HD-N11.

We have proved an injection of the quotient in (N16.5) into
$\ker\sigma^2_{E_\delta}$, and consequently
\begin{equation}
 \dim\ker\sigma^2_{E_\delta}\ \geq\
  \operatorname{ext}^3(P,N)
  -\operatorname{ext}^1(P,P)-\operatorname{ext}^1(N,N). \tag{N16.7}
\end{equation}
This does not assume a rank for any Yoneda matrix involving
$\delta$; it bounds the dimension of its exact quotient for
every such matrix.

\paragraph{3. Exact Ext dimensions for the specified sheaves.}
Here is a short finite formula determining every dimension used
in (N16.7). For a single elliptic pair, use slopes
$a,b\in\{-1,0,1\}$ and let $r$ and $t$ be the source and target
twist exponents, each either $0$ or $1$. Put $n_a=d+a^2$.
The Ext-dimension generating polynomial on that pair is
\begin{equation}
 \mathcal E(a,r;b,t;z)=
 \begin{cases}
   (a-b)^2z,&a\ne b,\\
   1+2z+z^2,&a=b,\ t=r,\\
   n_a(1+z),&a=b,\ t-r=1,\\
   n_a(z+z^2),&a=b,\ t-r=-1.
 \end{cases}                                        \tag{N16.8}
\end{equation}
For distinct slopes, the two graph curves meet transversely at
$(a-b)^2$ reduced points. Their local Koszul resolution has only
degree-one Ext, of length one at each intersection; line twists
do not change those dimensions. For equal slopes, the normal
line is trivial. The codimension-one local Ext sequence has
the line bundle on the elliptic curve in sheaf Ext degrees
zero and one. Its local-to-global spectral sequence has no
nonzero differential because the support has dimension one.
Elliptic Riemann--Roch gives exactly the other three polynomials.
For equal twists the line ratio is literally trivial, since the
same $M$ was used throughout. K\"unneth gives the product of
these four polynomials on $A$.

Let $\mathcal S_+$ and $\mathcal S_-$ be the package--slope pairs
with respectively positive and negative $c_\epsilon$. For
$U=(J,\epsilon)$ write $m_U=|c_\epsilon|$ and
$r_j(U)=1_{j\notin J}$. The exact finite sums are
\begin{align}
 A(d)&=\sum_{U,V\in\mathcal S_+}m_Um_V[z]
       \prod_{j=1}^4\mathcal E(\epsilon_j,r_j(U);
                                      \eta_j,r_j(V);z),\notag\\
 B(d)&=\sum_{U\in\mathcal S_+,V\in\mathcal S_-}m_Um_V[z^3]
       \prod_{j=1}^4\mathcal E(\epsilon_j,r_j(U);
                                      \eta_j,r_j(V);z),\notag\\
 C(d)&=\sum_{(J,\epsilon)\in\mathcal S_-}
                      |c_\epsilon|\prod_{j\in J}(d+\epsilon_j^2).
                                                        \tag{N16.9}
\end{align}
The involution $\epsilon\mapsto-\epsilon$ proves that the
negative-negative degree-one sum also equals $A(d)$.
Equation (N16.3) gives
$\operatorname{ext}^3(L,N)=2C(d)$ and
$\operatorname{ext}^2(L,N)=C(d)$.
It also gives $\operatorname{Ext}^1(L,P_0)=0$ and
$\operatorname{Ext}^7(L,P_0)=0$. Serre duality on $A$, whose
canonical bundle is trivial, then gives
$\operatorname{Ext}^1(P_0,L)=0$. Since
$\operatorname{ext}^1(L,L)=8$, the dimensions in (N16.7) are
\begin{equation}
 \operatorname{ext}^1(P,P)=A(d)+8,\qquad
 \operatorname{ext}^1(N,N)=A(d),\qquad
 \operatorname{ext}^3(P,N)=B(d)+2C(d).                 \tag{N16.10}
\end{equation}

Expanding the explicitly finite sums (N16.9) gives
\begin{align*}
 A(d)&=768\bigl(27d^8+146d^7+394d^6+644d^5+713d^4
                   +554d^3+298d^2+100d+16\bigr),\\
 B(d)&=192(d+1)\bigl(304d^8+3537d^7+17733d^6+44974d^5
                   +58594d^4\\
     &\hspace{37mm}{}+44689d^3+20089d^2+5168d+640\bigr),\\
 C(d)&=24d(d+1)^2(35d^2+35d+8).
\end{align*}
For clarity, absolute values in (N16.9) introduce no unstated
piecewise choice. Substitute $d=1+t$ into the coefficient formula
(N16.1). The recurrence
\[
 (R,S)\longmapsto\bigl(-dR-d\epsilon_jS,
                                      \epsilon_jR-dS\bigr)
\]
starts at $(1,0)$, and the resulting polynomial $c_\epsilon(1+t)$
has either only nonnegative coefficients or only nonpositive
coefficients, for each of the $81$ grid points. This can be
checked by that four-step recurrence; the complete exact
polynomial certificate below checks all $81$ instances.
Zero constant terms cause no problem at $d=1$: the corresponding
multiplicity is then zero and contributes zero to each sum.

The difference $B(d)+2C(d)-2A(d)-8$, written in $t=d-1$, is
\begin{align}
 S(t)={}&70\,732\,408+315\,691\,968t+614\,745\,312t^2
                 +683\,183\,760t^3\notag\\
       &+475\,906\,176t^4+214\,886\,544t^5
                 +62\,842\,560t^6\notag\\
       &+11\,528\,832t^7+1\,221\,312t^8+58\,368t^9.
                                                        \tag{N16.11}
\end{align}
Every displayed coefficient is positive. Equations (N16.7),
(N16.10) and (N16.11) prove the stated uniform lower bound for
every integer $d\geq1$ and every $\delta$.

\paragraph{4. Point stabilizers and the precise inherited actions.}
The standard cohomology sheaves of (N16.2) are exactly
$H^{-1}(E_\delta)=N$ and $H^0(E_\delta)=L\oplus P_0$.
The torsion subsheaf of $H^0$ is intrinsically $P_0$, and its
torsion-free quotient is intrinsically $L$. Suppose a point
functor $t_a^*(-)\otimes Q$, $Q\in\operatorname{Pic}^0(A)$,
preserves $E_\delta$. It preserves these sheaves and their
support components. Translation preserves the tangent space
of a graph, so a component $B_\epsilon$ can only be carried to
the same graph slope. Every $B_{e_j}$ occurs positively by
(N16.1), whence
\[
                a\in\bigcap_{j=1}^4B_{e_j}=\{0\}.
\]
Indeed the equations $z'_k=\delta_{jk}z_k$ for every $j,k$
force every $z_k$ and $z'_k$ to vanish. The rank-one quotient
then gives $L\otimes Q\simeq L$, hence $Q\simeq\mathcal O_A$.
The original point stabilizer is therefore trivial.

An actual Fourier--Mukai equivalence identifies self-Ext and
categorical traces and carries point functors by its actual
graph isomorphism, as established in HD-G19 from the full
GLO Proposition4.3.7 proof and Orlov Theorem2.10. Thus the
nonzero kernel and the original trivial stabilizer transport
along that precise equivalence.

Let $p:A'\to A$ be a finite \emph{\'etale} isogeny. The adjunction
map followed by the finite trace, divided by $\deg p$, gives
\[
 E_\delta\longrightarrow Rp_*p^*E_\delta
                         \longrightarrow E_\delta
                         \quad\text{equal to the identity}.
\]
Equivalently $\mathcal O_A$ is the split constant summand of
$p_*\mathcal O_{A'}$. Projection formula and adjunction then
show that pullback is injective on all self-Ext groups.
Its image is fixed by the canonical coherent deck action on
$p^*E_\delta$. Atiyah trace is natural for the \emph{\'etale}
pullback, so the nonzero classes in (N16.4)--(N16.7) remain
trace-zero and deck-invariant. For the full deck group, actual
finite \emph{\'etale} descent identifies this equivariant object
with $E_\delta$ on $A$. For a subgroup $H$ of the deck group,
the descended object on $A'/H$ is literally $p_H^*E_\delta$,
where $p_H:A'/H\to A$ is the induced finite \emph{\'etale}
map. Pullback to this quotient is again injective on self-Ext
and preserves the zero trace. Descent therefore identifies a
nonzero kernel of the \emph{equivariant} categorical trace;
ordinary deck invariance alone is not being used to infer that
last assertion. An actual equivalence transports object, action,
Ext and trace together.
This proves precisely the inherited-action assertion made at
the start. A different post-isogeny coherent action requires
its own invariant quotient calculation; none is inferred
from the ordinary dimension surplus.

\paragraph{5. Exact certificates and review scope.}
The dependency-free script
\texttt{six\_\allowbreak package\_\allowbreak ext\_\allowbreak polynomial.py}
reconstructs the
coefficients in $\mathbf Z[t,s]/(s^2+t+1)$, verifies their
uniform signs, multiplies (N16.8), and evaluates all sums in
the polynomial ring $\mathbf Z[t]$. Its exact JSON output
records (N16.10)--(N16.11). The separate script
\texttt{six\_\allowbreak package\_\allowbreak ext\_\allowbreak dimensions.py}
directly sums the
integer formulas for
$d=1,2,3,7,11,19,43,67,163$; all nine values agree with the
polynomial certificate. The smallest and largest values in
that list are respectively $70\,732\,408$ and
$5\,099\,767\,445\,971\,078\,738\,387\,960$.
An independent reconstruction by \texttt{/root/g34}, without
importing these scripts or their data, agrees with all nine
integer sums. These checks concern the finite algebra; the
exact quotient injection, trace identity and inherited-action
argument are proved above and are recorded separately from
computational checks in the review log.
